\documentclass[12pt]{article}
\usepackage[margin=1in]{geometry}
\title{Document with Footnotes}
\author{Footnote Tester}
\begin{document}
\maketitle

\section{Introduction}
Footnotes are commonly used in academic writing\footnote{This is a simple
footnote that provides additional context.} to provide supplementary information
without cluttering the main text.

\section{Technical Details}
The PDF format stores footnotes as regular text positioned at the bottom of the
page\footnote{PDF does not have a semantic concept of ``footnotes'' --- they are
simply text blocks positioned near the page bottom.}. This means the converter
must use spatial analysis to identify them.

Our approach detects footnotes by looking for small-font text blocks near the
bottom margin\footnote{Specifically, we check for text blocks where the font
size is less than 80\% of the body font size and the vertical position is in
the bottom 15\% of the page.}. These blocks are then tagged as footnotes in the
intermediate representation.

\section{Limitations}
Multi-page footnotes\footnote{Footnotes that are too long to fit on a single
page will overflow to the next page. This is a known edge case that our current
implementation does not handle perfectly.} present additional challenges.

\section{Conclusion}
Footnote handling is important for academic document conversion. The key signals
are: small font size, bottom-of-page positioning, and superscript reference markers
in the body text.
\end{document}